\begin{savequote}[8cm]
Alles Gescheite ist schon gedacht worden.\\
Man muss nur versuchen, es noch einmal zu denken.

All intelligent thoughts have already been thought;\\
what is necessary is only to try to think them again.
  \qauthor{--- Johann Wolfgang von Goethe \cite{von_goethe_wilhelm_1829}}
\end{savequote}

\minitoc

\section{Abstract}
The phenomenon of neutrinos changing from one flavour into another as they propagate is referred to as neutrino oscillations. The long-baseline neutrino oscillation experiments \acrshort{dune} and \acrshort{t2k} aim to precisely measure the neutrino oscillation parameters describing this effect. However, in addition to measuring the standard ($3$$\times$$3$-\acrshort{pmns}) neutrino oscillation parameters, accelerator-neutrino beams can be used to study new Beyond-the-Standard-Model physics models such as quantum decoherence that can manifest in quantum gravitational effects, \acrlong{liv} or \acrlong{nsi}. For this study, I have implemented quantum decoherence in the oscillation analyses that are already used within the long-baseline experiments \acrshort{dune} and \acrshort{t2k}. Moreover, I have studied the sensitivity to the parameters underpinning the decoherence model via sample variations. The goal of this project is to determine the sensitivity of \acrshort{dune} and \acrshort{t2k} to quantum decoherence signatures and set upper limits on the strength of this effect.